\section{Protocolo de pruebas}
\label{sec:pruebas}

El protocolo cubre los tipos de prueba de un sistema informático, desde las piezas
aisladas hasta el sistema completo bajo carga y fallos. Las pruebas se ejecutan en
este orden: si un nivel falla, no se avanza al siguiente.

\subsection{Tipos de prueba, herramientas y criterios}

\begin{longtable}{p{3.0cm}p{5.6cm}p{3.2cm}p{3.5cm}}
\caption{Tipos de prueba.}\\
\toprule
Tipo & Qué verifica & Herramientas & Cuándo \\
\midrule
\endfirsthead
\toprule
Tipo & Qué verifica & Herramientas & Cuándo \\
\midrule
\endhead
Análisis estático & Errores de programación sin ejecutar: variables sin usar,
conversiones peligrosas, estilo. & \codigo{g++ -Wall -Wextra}, \codigo{cppcheck},
\codigo{clang-tidy} & En cada compilación \\
Unitarias & Cada clase por separado: \codigo{Mensaje} (serializar y deserializar),
\codigo{FechaEvento} (\codigo{hayCupo}), \codigo{BaseDatos} (reglas de reserva),
\codigo{AlmacenamientoArchivo} (bitácora). & GoogleTest & En cada cambio \\
Integración & Dos componentes reales juntos: RMC $\leftrightarrow$ persistencia,
gestor $\leftrightarrow$ RMC, reconstrucción desde la bitácora. & GoogleTest + procesos
reales en una máquina & En cada cambio \\
Protocolo (contrato) & Cada proceso envía y acepta exactamente las tramas y campos
acordados; errores ante mensajes mal formados. & Cliente de prueba con cppzmq & En
cada cambio \\
Funcionales (extremo a extremo) & Cada requisito funcional a través de las tres
capas (Tabla~\ref{tab:e2e}). & \codigo{bin/cliente} con guiones de entrada + script
\codigo{bash} & Antes de cada entrega \\
Concurrencia y consistencia & Sin sobrecupo, sin reservas duplicadas, contadores
iguales a la suma de reservas. & Varios clientes simultáneos & Antes de cada entrega \\
Regresión & Lo que funcionaba sigue funcionando después de un cambio. & Repetir las
suites anteriores & En cada \emph{merge} a \codigo{develop} \\
Desempeño (Tabla~1) & Tiempo promedio, desviación estándar y TPM de consultas y
reservas con 3, 5 y 8 clientes, base frente a escalable. & Sección~\ref{sec:metricas}
y \ref{sec:carga} & Entrega final \\
Carga, estrés y resistencia & Meta de p95 < 1\,s, punto de saturación y
estabilidad durante horas (fugas de memoria). & Generador de carga, \codigo{pidstat},
Valgrind & Entrega final \\
Escalabilidad & Ganancia de TPM al pasar de 1 a 2 y 3 procesadores RMC. &
Generador de carga & Entrega final \\
Inyección de fallos & Caídas y mensajes perdidos se enmascaran
(Tabla~\ref{tab:fallos}). & \codigo{kill}, \codigo{tc netem}, \codigo{iptables} &
Entrega final \\
Seguridad & Validación de entradas, propiedad de las reservas, repetición de
mensajes, tráfico cifrado. & Cliente de prueba, Wireshark & Entrega final \\
Despliegue e instalación & El sistema se compila y arranca en máquinas limpias
siguiendo el README. & \codigo{make}, scripts, Kubernetes & Antes de cada entrega \\
Aceptación y usabilidad & Un estudiante completa consulta y reserva sin ayuda; los
mensajes son claros. & Sesión guiada con usuarios & Entrega final \\
\bottomrule
\end{longtable}

\subsection{Casos funcionales}

\begin{longtable}{p{1.4cm}p{6.6cm}p{7.3cm}}
\caption{Casos de prueba funcionales de extremo a extremo.}
\label{tab:e2e}\\
\toprule
ID & Caso & Resultado esperado \\
\midrule
\endfirsthead
\toprule
ID & Caso & Resultado esperado \\
\midrule
\endhead
F-01 & Consultar un mes con eventos (\codigo{2026-11}). & Solo fechas de ese mes, en
orden, con lugar, precio, límite, reservas y disponibles. \\
F-02 & Consultar un mes sin eventos. & Lista vacía, sin error. \\
F-03 & Consultar con un mes mal escrito (\codigo{2026-13}). & \codigo{DATOS\_INVALIDOS}. \\
F-04 & Reservar 2 cupos en una fecha con cupo. & Reserva creada; los disponibles
bajan en 2; se registra una transacción. \\
F-05 & Reservar más cupos de los disponibles. & \codigo{CUPOS\_INSUFICIENTES}; nada
cambia. \\
F-06 & Reservar exactamente los cupos restantes y luego uno más. & La primera
funciona; la segunda \codigo{CUPOS\_INSUFICIENTES}. \\
F-07 & Segunda reserva del mismo estudiante en el mismo evento. &
\codigo{RESERVA\_DUPLICADA}. \\
F-08 & Reservar con cupos 0 o negativos, o una fecha inexistente. &
\codigo{DATOS\_INVALIDOS} / \codigo{FECHA\_NO\_EXISTE}. \\
F-09 & Reservar con un estudiante no registrado. & \codigo{ESTUDIANTE\_NO\_REGISTRADO}. \\
F-10 & Modificar cupos y luego fecha de una reserva. & Contadores de ambas fechas
correctos; diferencia de precio calculada. \\
F-11 & Modificar o cancelar la reserva de otro estudiante. & \codigo{RESERVA\_AJENA}. \\
F-12 & Cancelar una reserva. & Respuesta inmediata \codigo{EN\_PROCESO}; después
notificación \codigo{CONFIRMADA} y cupos liberados. \\
F-13 & Modificar una reserva con cancelación pendiente. & \codigo{CANCELACION\_PENDIENTE}. \\
F-14 & Reiniciar persistencia y consultar. & Las reservas anteriores siguen ahí. \\
\bottomrule
\end{longtable}

\subsection{Concurrencia y fallos}

\begin{longtable}{p{1.4cm}p{6.6cm}p{7.3cm}}
\caption{Casos de concurrencia e inyección de fallos.}
\label{tab:fallos}\\
\toprule
ID & Caso & Resultado esperado \\
\midrule
\endfirsthead
\toprule
ID & Caso & Resultado esperado \\
\midrule
\endhead
C-01 & 20 clientes piden 1 cupo a la vez en una fecha con 5 disponibles. & Exactamente 5
reservas; 15 \codigo{CUPOS\_INSUFICIENTES}. \\
C-02 & El mismo \codigo{idSolicitud} enviado 10 veces en paralelo. & Una sola
reserva; 10 respuestas idénticas. \\
C-03 & Al final de cada corrida: suma de cupos de las reservas activas de cada
fecha. & Igual al contador de reservas de la fecha y nunca mayor al límite. \\
FT-01 & Matar un procesador RMC durante una ráfaga de reservas. & Los clientes
reintentan; sin duplicados; el gestor lo descarta en $\le$ 3\,s. \\
FT-02 & Perder la respuesta de \codigo{reservar} (retardo con \codigo{tc netem}
mayor al timeout). & El reintento devuelve la misma reserva, sin duplicado. \\
FT-03 & Reiniciar el gestor con carga. & Los procesadores se vuelven a registrar;
los clientes continúan tras reintentar. \\
FT-04 & Detener el servicio de cancelación y cancelar reservas. & Las solicitudes
quedan \codigo{PENDIENTE}; al volver, el barrido las aplica una vez. \\
FT-05 & Matar persistencia a mitad de la carga y reiniciarla. & Ninguna reserva
aplicada a medias; el estado se reconstruye desde la bitácora. \\
FT-06 & 10\,\% de pérdida de paquetes entre cliente y gestor. & Los reintentos
funcionan; sin duplicados. \\
\bottomrule
\end{longtable}

\subsection{Pruebas de desempeño}

\begin{itemize}
    \item \textbf{Escenarios de la Tabla~1} (sección~\ref{sec:metricas}): 3, 5 y 8
          clientes en las arquitecturas base y escalable.
    \item \textbf{Carga:} subir gradualmente el número de clientes hasta el uso esperado
          y verificar p95 < 1\,s.
    \item \textbf{Estrés:} seguir aumentando clientes hasta que el p95 supere 1\,s o los
          errores superen 1\,\%; se registra el TPM máximo y qué componente se satura.
    \item \textbf{Resistencia:} carga media durante 2 horas; la memoria de los procesos
          no debe crecer sin límite.
    \item \textbf{Escalabilidad:} misma carga con 1, 2 y 3 procesadores RMC; se espera
          que el TPM crezca al agregar réplicas.
\end{itemize}

\subsection{Criterios de aceptación}

\begin{itemize}
    \item Todos los casos funcionales F-01 a F-14 pasan.
    \item Ninguna corrida termina con sobrecupo, reservas duplicadas o contadores
          descuadrados (C-03).
    \item El 95\,\% de las transacciones termina en menos de 1\,s en los escenarios de
          la Tabla~1.
    \item Los fallos FT-01 a FT-06 no producen pérdida ni duplicación de datos.
\end{itemize}

\subsection{Trazabilidad}

\begin{table}[H]
\centering
\caption{Requisitos frente a pruebas.}
\begin{tabularx}{\textwidth}{p{6.4cm}Y}
\toprule
Requisito & Pruebas \\
\midrule
Consultar eventos de un mes & Unitarias de \codigo{BaseDatos}; F-01 a F-03 \\
Reservar sin superar el límite & F-04 a F-09; C-01, C-03 \\
Modificar y cancelar (asíncrona) & F-10 a F-13; FT-04 \\
Reintentos sin duplicados & C-02; FT-01, FT-02, FT-06 \\
Atomicidad y durabilidad & F-14; FT-05 \\
Transacciones < 1\,s y Tabla~1 & Pruebas de desempeño \\
Seguridad & F-11; pruebas de seguridad \\
\bottomrule
\end{tabularx}
\end{table}
